\chapter{Introducere}
\section{Contextul general și motivația cercetării}
\hspace{0.75cm}Evoluția omului a constat în dezvoltarea gândirii critice, structurate, logice cu scopul de a depăși 
impedimentele puse în calea existentei acestuia. 
Cu cât sfera cunoasterii s-a largit, indivizii au fost nevoiți să se specializeze în anumite domenii,
cum ar fi agricultura, astronomia, arhitectura sau medicina. 
Pentru a putea continua evolutia acestora, apare necesitatea individului cu o 
gândire pretata pentru probleme complexe interdisciplinare: anume, inginerul. 
Urmărind avansul cunoștințelor din domeniul ingineriei, această 
meserie s-a ramificat în diferite directii, precum: mecanica, energetica, chimica, electronică, telecomunicații, 
automatică, etc.

\hspace{0.75cm}Automatica reprezintă o știință interdisciplinară care are o viziune holistică asupra sistemelor complexe. 
Scopul acesteia este de a optimiza procesele si evolutia acestor sisteme, minimizand contribuția si eroarea umana.
Unul dintre cele mai complexe sisteme cunoscute de umanitate este reprezentat de insusi corpul omenesc. 

\hspace{0.75cm}O perturbatie in urma careia acest sistem rareori revine la un regim de funcționare normal  
este accidentul vascular cerebral. La nivel global, AVC-ul reprezintă o problemă majoră de sănătate publică, 
fiind a doua cauză de mortalitate și a treia cauză principală de dizabilitate pe termen lung, conform studiului 
Global Burden of Disease \cite{gbd2021}. Efectele unui AVC sunt profunde și adesea ireversibile, 
privând rapid porțiuni ale creierului de oxigen și nutrienți,
ceea ce duce la necroză tisulară și la pierderea unor funcții motorii, cognitive sau senzoriale esențiale.

\hspace{0.75cm}În România, situatia este critică, incidența și gravitatea acestui fenomen atingând cote semnificativ
mai mari  în comparație cu media europeană. Statisticile plasează România în topul țărilor europene în ceea ce 
privește mortalitatea cauzată de afecțiuni cerebrovasculare, rata decesului prin AVC fiind de aproximativ două ori
mai mare decât media Uniunii Europene și reprezentând peste 20\% din totalul deceselor la nivel național 
\cite{romania_stats}. Anual, aproximativ 60 de mii de pacienți români suferă un accident vascular cerebral ischemic sau hemoragic. 

\hspace{0.75cm}Un alt capitol la care Romania se afla in coada clasamentelor europene este reprezentat de accesul facil la 
infrastructura medicală și tratamente de recuperare. Majoritatea centrelor de recuperare neurologică bine echipate 
și a specialiștilor cu experiență în neuroreabilitare sunt concentrate în marile centre universitare \cite{tiu_2023}.

\section{Obiectivele proiectului}

\hspace{0.75cm} Luând în considerare problemele menționate anterior, în cadrul acestei lucrări este abordată 
următoarea soluție: realizarea unui sistem telemedical integrat, compus dintr-o aplicație de tip wearable (pentru smartwatch) 
și o aplicație mobilă, care să monitorizeze, să cuantifice și să personalizeze de la distanță procesul de recuperare motorie 
a pacienților, procesul de reabilitare fiind asistat de algoritmi de procesare de semnal și modele de tip machine learning.

\hspace{0.75cm}Obiectivele specifice ale proiectului sunt: 
\begin{itemize}
    \item Dezvoltarea unei aplicații multiplatform, capabilă să monitorizeze în 
     timp real mișcările și semnalele fiziologice ale pacienților.
    \item Integrarea algoritmilor specifici ingineriei biomedicale în vederea evaluării și cuantificării 
     numerice obiective a calității mișcărilor efectuate de pacient. 
    \item Integrarea algoritmilor de machine learning astfel încât sistemul să poată calcula un progres realist 
     și să ajusteze automat obiectivele de recuperare ale pacientului pe diferite orizonturi de timp
     (zilnic, săptămânal, lunar).
    \item Oferirea unui mod de vizualizare intuitiv care să permită o metodă ușoară de urmărire a 
     progresului și de comparare a sesiunilor de recuperare, atât pentru pacient, cât și pentru personalul medical. 
    \item Asigurarea securității și confidențialității datelor medicale colectate  
\end{itemize}
